
\section{Round-twenty-seven reconstruction: causal subtree shears and a complete collision-history algebra}
\label{sec:r27-b2}

This paper supplies the hard-sphere root input.  The principal new device is a
causal subtree-shear coordinate for every surplus contact.  It gives a
block-triangular loop Jacobian with a nonzero diagonal on the regular set, so
cyclomatic number is no longer confused with geometric rank and no analytic
determinant ideal is allowed to vanish identically.  The same causal
decomposition defines an associative cut-history algebra and a factorial
Duhamel estimate.

\subsection{Microscopic state and contact measure}

Let $N$ hard spheres of diameter $\varepsilon$ move on $\mathbb T^3$ under
elastic reflection, with Boltzmann--Grad activity
$\mu_\varepsilon\varepsilon^2=1$.  A regular contact mark is
\[
 c=(t,x,v,v_*,\omega),\qquad
 B(c)=[(v-v_*)\cdot\omega]_+.
\]
For a density $f$ define
\[
 A_f(dc)=\tfrac12 f(t,x,v)f(t,x,v_*)B(c)
 \,dt\,dx\,dv\,dv_*\,d\omega.                             \tag{B2.1}
\]
The regular preparation class satisfies a Gaussian upper bound, a positive
Maxwellian lower bound on compact velocity sets, and
\[
 \sup_{t,x}\|f(t,x,\cdot)\|_{L^\infty}
 +\sup_{t,x}\int e^{a|v|^2}f(t,x,v)\,dv<\infty.            \tag{B2.2}
\]

\begin{lemma}[Uniform relative-velocity anti-concentration]
\label{lem:r27-b2-relative}
On (B2.2),
\[
 \int_{\mathbb R^3\times\mathbb R^3}
 {f(v)f(v_*)\over|v-v_*|}\,dv\,dv_*\le C,                 \tag{B2.3}
\]
and hence
\[
 A_f\{0<(v-v_*)\cdot\omega<\eta\}\le C\eta^2.             \tag{B2.4}
\]
If $\Gamma=qA_f$ and $\|q\|_{L^p(A_f)}\le L$, then
$\Gamma(G_\eta)\le C_L\eta^{2(1-1/p)}$.
\end{lemma}

\begin{proof}
Set $g=v-v_*$.  The convolution
$\rho(g)=\int f(v)f(v-g)\,dv$ is bounded by
$\|f\|_\infty$ and has Gaussian tail.  Since $|g|^{-1}$ is locally integrable
in three dimensions, $\int\rho(g)|g|^{-1}dg<\infty$.  Direct angular
integration gives
\[
 \int_{0<g\cdot\omega<\eta}[g\cdot\omega]_+\,d\omega
 \le C\eta^2/|g|
\]
for $|g|>\eta$, while the region $|g|\le\eta$ is $O(\eta^4)$ by boundedness of
$\rho$.  This proves (B2.4); Holder gives the last claim.  Thus the constant
does not rely on a moment bound alone.
\end{proof}

\subsection{Decorated histories and the cut product}

A decorated collision history
\[
 G=(V,E,\prec,\ell,\mathbf z,\mathbf o)
\]
consists of a finite rooted directed forest, a strict chronological order on
vertices and collision edges, particle labels $\ell$, continuous collision
data $\mathbf z$, and a finite ordered list $\mathbf o$ of open half-edges
crossing a time cut.  Every edge carries its pre- and post-collisional
velocities and the exact elastic involution.  Let $\Omega_G$ be the
semianalytic parameter manifold obtained by imposing all tree contacts in
chronological order and removing grazing and simultaneous contacts.

For $a>0$ define
\[
 \mathfrak H_a
 =\left\{\mathbf H=(H_G):
 \|\mathbf H\|_a
 =\sum_G {a^{|V(G)|}\over|\operatorname{Aut}G|}
 \int_{\Omega_G}e^{\alpha\sum_{\ell}|v_\ell|^2}\,d|H_G|<\infty\right\}.
                                                                    \tag{B2.5}
\]
The automorphism factor equals the exponential-generating-function
normalization; for fully labelled histories it is $|V|!$.

If $\mathbf H^-$ and $\mathbf H^+$ have matching open half-edges, define
$\mathbf H^-\star\mathbf H^+$ by the fibre product over their ordered labels,
integrating the common cut state once and summing over compatible label
partitions.

\begin{proposition}[Associative collision-history algebra]
\label{prop:r27-b2-algebra}
For $0<a'<a$,
\[
 \|\mathbf H^-\star\mathbf H^+\|_{a'}
 \le C(a,a')\|\mathbf H^-\|_a\|\mathbf H^+\|_a.            \tag{B2.6}
\]
The product is measurable, positive on positive histories, and associative.
Density and contact current are continuous linear projections of
$\mathfrak H_a$.
\end{proposition}

\begin{proof}
For a fixed glued graph, the ordered label partitions are precisely the
orbits of the product automorphism group.  The exponential generating factors
therefore satisfy the binomial convolution identity.  Velocity weights are
submultiplicative under elastic reflection because kinetic energy is
preserved.  The loss from summing open-half-edge matchings is bounded by
$(a-a')^{-r}$ and is absorbed by the radius loss.  Fubini on the common cut
variables shows that the two possible parenthesizations are the same triple
fibre product, proving associativity.  Positivity and continuity of the
projections are immediate from the definition.
\end{proof}

\subsection{Causal subtree-shear coordinates}

Fix a connected ordered graph $G$ and a chronological spanning tree $T$.
Each surplus edge $e\in E\setminus T$ closes a loop at time $t_e$.  Cut the
unique tree edge immediately preceding the later endpoint of $e$ and let
$D_e$ be the descendant subtree.  Before $t_e$, introduce a shear parameter
$\zeta_e\in\mathbb R^3$ by translating the root position of $D_e$ by
$\zeta_e$ and transporting every later collision in $D_e$ by the same
translation until it first meets the complement.  Earlier imposed contacts
inside $D_e$ and its complement are unchanged.

Let
\[
 \mathcal C_e(\zeta)
 =x_{i_e}(t_e^-;\zeta)-x_{j_e}(t_e^-;\zeta)-\varepsilon\omega_e.         \tag{B2.7}
\]
Order surplus edges chronologically.

\begin{lemma}[Causal triangular loop Jacobian]
\label{lem:r27-b2-triangular}
On every regular chronology chart,
\[
 {\partial\mathcal C_e\over\partial\zeta_{e'}}
 =0\quad\text{if }t_{e'}>t_e,\qquad
 {\partial\mathcal C_e\over\partial\zeta_e}
 =\pm I+R_e,                                             \tag{B2.8}
\]
where $\|R_e\|\le C\delta_{\rm sing}$ after excluding a
$\delta_{\rm sing}$-neighbourhood of grazing, simultaneous, and tangential
tree--complement contacts.  Thus, for $\delta_{\rm sing}$ fixed small,
\[
 \left|\det{\partial(\mathcal C_e)_{e\in E\setminus T}
                \over\partial(\zeta_e)_{e\in E\setminus T}}\right|
 \ge2^{-3s(G)},\qquad s(G)=|E|-|V|+1.                    \tag{B2.9}
\]
\end{lemma}

\begin{proof}
A shear introduced for a later surplus edge does not change a trajectory at
an earlier time, giving the zero block.  At time $t_e$, exactly one endpoint
of the closing edge belongs to $D_e$, so direct translation contributes
$\pm I$.  A prior collision of the translated subtree with its complement
would change the derivative; by construction the first such contact is the
closing event, except on the excluded chronology/simultaneity set.  Near a
regular reflection, derivative transport differs from the identity by a
quantity controlled by the distance to that exceptional set.  Hence the
Jacobian is block lower triangular with uniformly invertible diagonal blocks.
This is a geometric rank theorem, not an application of Lojasiewicz to an
ideal that might vanish identically.
\end{proof}

\begin{theorem}[Multiplicative surplus-contact gain]
\label{thm:r27-b2-loop}
For every fixed vertex cutoff $K$ and regular tilted class,
\[
 |\mathcal I_G^\varepsilon|
 \le C_K^{|E(G)|}\varepsilon^{\alpha_Ks(G)}
      |\mathcal I_{T_G}^\varepsilon|                     \tag{B2.10}
\]
with $\alpha_K>0$.  The gains multiply for all surplus edges.
\end{theorem}

\begin{proof}
On the regular set, coarea in all shear variables uses (B2.9) and gives one
tube factor $C\varepsilon$ per independent scalar closure coordinate; after
the ordinary cross-section/activity cancellation, a positive power remains
for every surplus edge.  The excluded set is covered by grazing,
simultaneous-contact, and tangential charts.  Lemma
\ref{lem:r27-b2-relative} and elementary time/angular coarea bound it by
$C\delta_{\rm sing}^{\kappa_K}$.  Optimize
$\delta_{\rm sing}=\varepsilon^{\theta_K}$.  Because the full loop Jacobian is
triangular, all closure variables are integrated simultaneously and their
gains multiply without recomputing an unverified determinant ideal.
\end{proof}

\subsection{The factorial Duhamel map}

Let $\mathcal T_\Delta^\varepsilon$ be the exact history Duhamel map on a time
slice of length $\Delta$.  Free transport preserves (B2.5).  A creation edge
adds one vertex and one time integral; cutting a history at the slice boundary
retains all ordered open half-edges.

\begin{lemma}[Radius-loss Duhamel estimate]
\label{lem:r27-b2-duhamel}
For $0<a'<a$,
\[
 \|\mathcal T_\Delta^\varepsilon(\mathbf H)
      -\mathcal T_\Delta^\varepsilon(\widetilde{\mathbf H})\|_{a'}
 \le {C\Delta\over a-a'}
 \exp\!\left({C\Delta\over a-a'}\right)
 \|\mathbf H-\widetilde{\mathbf H}\|_a,                   \tag{B2.11}
\]
uniformly in $\varepsilon$.  Histories with more than $K$ new vertices have
norm at most
\[
 C\exp(C\Delta/(a-a'))(a'/a)^K.                           \tag{B2.12}
\]
\end{lemma}

\begin{proof}
Differentiate the exponential generating series with respect to one newly
created labelled vertex.  The binomial identity behind Proposition
\ref{prop:r27-b2-algebra} turns the sum of insertion sites into
$\partial_a$, bounded by $(a-a')^{-1}$ through Cauchy's estimate.  Iterating
the time-ordered integral gives $\Delta^k/k!$ and sums to the exponential in
(B2.11).  Marking $K$ vertices before evaluating at $a'$ gives (B2.12).
\end{proof}

\begin{theorem}[Finite-time hierarchy propagation]
\label{thm:r27-b2-propagation}
For every regular Boltzmann time $T<T_*$ there is $a_T>0$ such that
\[
 \sup_{t\le T}\|\mathbf H_t^\varepsilon-\mathbf H_t\|_{a_T}\to0.         \tag{B2.13}
\]
The limiting hierarchy is the tree-history hierarchy and composition across
time cuts is exactly $\star$.
\end{theorem}

\begin{proof}
Partition $[0,T]$ into slices satisfying the radius-loss condition in Lemma
\ref{lem:r27-b2-duhamel}.  First cut at $K$; Theorem
\ref{thm:r27-b2-loop} removes all surplus histories in that finite family.
Then let $\varepsilon\to0$, and finally $K\to\infty$ using (B2.12).
Associativity of $\star$ makes the result independent of the partition.
\end{proof}

\subsection{Connected pressure}

For an analytic source $H$ define the connected coefficient by the logarithm
in the history algebra.  The tree-graph recursion and the automorphism factor
give
\[
 |C_k^\varepsilon(H)|\le C(CT)^k k^{-3/2},                \tag{B2.14}
\]
on a short slice; source derivatives mark finitely many vertices and preserve
summability.

\begin{theorem}[Grand-canonical analytic pressure]
\label{thm:r27-b2-pressure}
For $T<T_*$,
\[
 \Lambda_\varepsilon(H)
 ={1\over\mu_\varepsilon}\log
 E_{\rm GC}^\varepsilon e^{\mu_\varepsilon\langle H,
 (f^\varepsilon,\Gamma^\varepsilon,\mathbf H^\varepsilon)\rangle}
\]
converges locally uniformly, with every fixed source derivative, to an
analytic functional $\Lambda_T(H)$.  Its Hamiltonian is the
Hamilton--Boltzmann collision exponential and all surplus-history
coefficients vanish.
\end{theorem}

\begin{proof}
On a short slice, (B2.14) gives normal convergence and Theorem
\ref{thm:r27-b2-loop} eliminates surplus graphs term by term under a summable
majorant.  Theorem \ref{thm:r27-b2-propagation} and the associative cut
product concatenate the logarithmic series.  Cauchy estimates give uniform
source derivatives.  The last tree edge separates two rooted subtrees, whose
exact recursion sums to the collision exponential.
\end{proof}

\subsection{Positive balance recovery without a circular graph inverse}

A finite velocity/cell partition has balance space
\[
 \mathcal Y_M=\{r:\langle r,1,v,|v|^2\rangle=0\}.
\]
Choose finitely many regular elastic collisions
$c_1,\ldots,c_J$ whose gain--loss vectors $g(c_j)$ span $\mathcal Y_M$.
A positive Maxwellian background assigns every $c_j$ mass at least $\delta_M$.

\begin{lemma}[Collision-simplex correction]
\label{lem:r27-b2-simplex}
Every sufficiently small $r\in\mathcal Y_M$ has coefficients
$\alpha_j(r)$ with
\[
 r=\sum_j\alpha_j(r)g(c_j),\qquad
 |\alpha(r)|\le C_M\|r\|,
\]
and the corrected collision multiplier remains positive after adding
$\sum_j\alpha_j(r)$ on small disjoint neighbourhoods of the $c_j$.
\end{lemma}

\begin{proof}
Select a basis of gain--loss vectors from the spanning family and use its
finite matrix right inverse.  Split positive and negative coefficients across
two copies of each collision neighbourhood.  The positive background
$\delta_M$ absorbs every sufficiently small correction.
\end{proof}

\begin{theorem}[Positive regular recovery and exact tilts]
\label{thm:r27-b2-recovery}
Every admissible path in the attainable action domain is approximated by
smooth positive pairs $(f_n,q_nA_{f_n})$ satisfying exact balance and
conservation, with convergent entropy action.  For each regular pair, the
exact deterministic finite-volume tilt concentrates exponentially at that
pair.
\end{theorem}

\begin{proof}
Truncate velocity and time, mollify $f$ and $q$, and add a vanishing
Maxwellian background.  The resulting balance defect lies in
$\mathcal Y_M$ and tends to zero.  Lemma \ref{lem:r27-b2-simplex} removes it
while preserving positivity.  Convexity of the perspective entropy and
uniform integrability give cost convergence.

Tilt the original deterministic initial law by the source dual to the regular
control.  The pressure theorem gives the normalizer.  Strict convexity of the
relative entropy on the positive chart and uniqueness of the controlled
Boltzmann equation make the tilted variational problem have one minimizer.
The second source derivative supplies an exponential variance bound; a finite
cover of the complement yields concentration.  No independent Poisson
collision process is introduced.
\end{proof}

\subsection{The correct LDP topology}

The microscopic density coordinate is the empirical probability measure,
with narrow convergence plus the declared velocity Orlicz moment.  Strong
local $L^1$ is used only for a mollified correlation density
$f^\varepsilon*\rho_{\delta(\varepsilon)}$, where
$\varepsilon\ll\delta(\varepsilon)\downarrow0$.  Contact current and the
history hierarchy are retained explicitly.

\begin{theorem}[Joint dynamic hard-sphere LDP]
\label{thm:r27-b2-ldp}
The laws of
\[
 (f^\varepsilon,\Gamma^\varepsilon,\mathbf H^\varepsilon)
\]
satisfy a good LDP at speed $\mu_\varepsilon$ on the narrow-Orlicz,
contact-weak, factorial-history topology.  On the attainable domain the rate
is
\[
 I_{\rm GC}(f,\Gamma)
 =H(f_0\mid f_0^{\rm ref})+\int\ell(q)\,dA_f,
 \qquad\Gamma=qA_f,                                      \tag{B2.15}
\]
subject to exact balance; outside the lower-semicontinuous attainable closure
it is $+\infty$.  B1 coefficient extraction gives the exact-number and
microcanonical contraction.
\end{theorem}

\begin{proof}
The pressure theorem gives the compact upper bound.  Exponential particle,
velocity, contact, and factorial-history estimates give exponential
tightness.  Theorem \ref{thm:r27-b2-recovery} supplies a rate-dense exposed
class and exact deterministic tilts give the lower bound.  Lower
semicontinuity identifies (B2.15) on the attainable domain.  B1 supplies a
source-uniform local denominator and coarea fibre, so division yields the
conditioned principle without changing the dynamic rate.
\end{proof}

The retained hierarchy, associative cut product, causal rank theorem,
factorial propagation, relative-velocity bound, positive recovery,
tilted concentration, and consistent topology are now explicit root
theorems rather than names or postulated inequalities.
